\documentclass[11pt]{article}
\usepackage[margin=1in]{geometry}
\usepackage[T1]{fontenc}
\usepackage{mathpazo}
\usepackage{amsmath,amssymb,amsthm}
\usepackage{booktabs,longtable,array,enumitem,microtype}
\usepackage[colorlinks=true,linkcolor=black,urlcolor=blue!55!black]{hyperref}

\newtheorem{invariant}{Invariant}
\newtheorem{requirement}{Requirement}
\newcommand{\MUST}{\textsc{must}}
\newcommand{\MUSTNOT}{\textsc{must not}}
\newcommand{\SHOULD}{\textsc{should}}
\newcommand{\MAY}{\textsc{may}}
\newcommand{\fold}{\operatorname{fold}}

\title{\textbf{Persistent World Layer Specification}\\[4pt]
\large Appearance, Operative State, Provenance, and Committed History}
\author{Flyxion\\\small Independent Researcher}
\date{September 2026\\Version 0.2 (adds \S\ref{sec:resonant}, an optional resonant-fold extension; all other sections unchanged from Version 0.1)}

\begin{document}
\maketitle

\begin{abstract}
This document converts the layered account of persistent worlds into an implementation-facing specification. It defines six distinct roles: appearance, operative state, committed history, observation, rendering, and semantic folding, while retaining provenance as a separate certified fold. It specifies authority boundaries, event and checkpoint envelopes, gate outcomes, deterministic fold requirements, replay claims, concurrency rules, failure handling, conformance tests, and language mappings for Rust and Go. The central constraint is that visual recurrence, operational continuity, and historical identity are different properties. Rendering is disposable and non-authoritative; operative state drives provisional evolution; committed history warrants identity at commitment boundaries; and current operative and provenance states are certified folds over the same causal event history. This is a normative engineering contract, not a claim that every interactive system requires the strongest persistence grade.
\end{abstract}

\tableofcontents

\section{Status, Scope, and Normative Language}

This specification applies to interactive or simulated systems that claim persistence across sessions, processes, repairs, software versions, or distributed authorities. It is deliberately independent of storage engine, network protocol, rendering technology, and simulation language. A conforming implementation can be local or distributed, symbolic or neural, discrete or continuous.

The words \MUST, \MUSTNOT, \SHOULD, and \MAY{} are normative. ``History'' means committed history unless explicitly qualified. ``State'' is never used without identifying whether it means operative state, provenance state, or a provisional state.

The specification does not require every numerical integration step, frame, sensor sample, or rejected packet to be stored forever. It requires explicit policies for deciding which facts become identity-bearing commitments, how retained summaries are certified, and how omissions or compression remain auditable.

\section{System Model}

At a certified causal frontier $\phi$, a world instance is represented by
\[
W_{\phi}=(H_{\phi},M_{\phi},Q_{\phi},C_{\phi}),
\]
where $H_{\phi}$ is the committed event history visible at the frontier, $M_{\phi}$ is the certified operative fold, $Q_{\phi}$ is the certified provenance and permission fold, and $C_{\phi}$ is checkpoint metadata, if a checkpoint exists. An observer receives an appearance
\[
x=R(M_{\phi},Q_{\phi},\omega,\rho),
\]
where $\omega$ specifies observer identity, viewpoint, channel, and authorization, and $\rho$ specifies disposable presentation conditions such as resolution, level of detail, sampling seed, compression, or accessibility format.

Between commitment boundaries, the simulation may evolve through provisional states
\[
\widetilde M(t+\Delta)=F_{\Delta}(\widetilde M(t);V,\Xi).
\]
Here $V$ identifies applicable rule, model, schema, and solver versions, while $\Xi$ identifies controlled stochastic inputs. A provisional state is not historically authoritative merely because it was used to compute the next step.

\subsection{Authority matrix}

\begin{longtable}{p{0.20\textwidth}p{0.30\textwidth}p{0.39\textwidth}}
\toprule
Object & Authoritative for & Not authoritative for \\
\midrule
Appearance $x$ & What an observer was presented under declared render conditions & Object identity, ownership, causal history, or an instantaneous global state \\
Provisional operative state $\widetilde M$ & Continuing computation inside an uncommitted interval & Certified persistence, provenance, or recovery after rollback \\
Certified operative state $M_{\phi}$ & Current dynamic queries and transitions at frontier $\phi$ & The warrant or lineage of its values \\
Provenance state $Q_{\phi}$ & Current ownership, permission, attribution, audit, and standing queries & Complete historical evidence unless the fold retains it explicitly \\
Committed history $H_{\phi}$ & Which commitments occurred, their causal relations, and which folds are warranted & Efficient per-frame or per-tick computation \\
Checkpoint $C_{\phi}$ & Accelerated reconstruction from a named frontier under declared versions and tolerances & Replacement of the event history it summarizes \\
\bottomrule
\end{longtable}

\section{Layer Contracts}

\subsection{Appearance}
\label{sec:appearance-section}

Appearance is an observer-relative projection returned through a presentation channel. It includes images, audio, maps, text descriptions, instrument traces, administrative views, or network view models. Two appearances may be equal while their source worlds differ in state, provenance, or history. Conversely, one certified world may yield many valid appearances.

An appearance record \MAY{} carry a content hash, frontier identifier, observer-policy version, renderer version, and presentation parameters. Those fields make the output reproducible or attestable; they do not make the output authoritative by themselves.

\begin{invariant}[Non-authority of appearance]
Equality of rendered outputs never entails equality of operative state, provenance state, or committed history.
\end{invariant}

\subsection{Operative state}

Operative state contains the smallest representation on which required transition and query operations are defined while preserving all distinctions those operations need. Typical members include identifiers, positions, velocities, fields, inventories, object relations, health values, environmental resources, and dynamically active constraints.

An implementation \MUST{} distinguish certified operative state from provisional operative state. A certified state names the history frontier and fold implementation that warrant it. A provisional state names its certified base frontier and the uncommitted computation that produced it. Rollback of provisional state \MUSTNOT{} alter committed history.

\begin{invariant}[Frontier binding]
Every certified $M$ is bound to exactly one declared causal frontier $\phi$, schema version, and physical-fold version. It is invalid to compare or replace certified states while omitting those bindings.
\end{invariant}

\subsection{Provenance and permission state}

The provenance state $Q$ is a current, queryable fold of historically active relations that affect meaning or admissibility without necessarily changing immediate dynamics. Examples include ownership, authorship, repair lineage, access rights, encumbrances, sanctions, disputed claims, evidence standing, and audit counters.

The separation between $M$ and $Q$ is logical, not necessarily physical. They may share a database or process. Their update functions \MUST{} remain independently specified. A cross-layer consequence, such as permission opening a physical door, \MUST{} occur through a typed admitted event or an explicitly versioned deterministic rule, never through an unrecorded side effect.

\subsection{Committed history}

History is an append-only set of immutable event envelopes with a causal order. A single-writer implementation \MAY{} use a sequence. A distributed implementation \SHOULD{} use a directed acyclic graph or another representation that preserves causal parents and concurrency without inventing a false wall-clock total order.

Appending a corrective event is permitted. Replacing, deleting, or semantically reinterpreting an existing event in place is forbidden. If legal or operational policy requires physical deletion, history \MUST{} retain an auditable tombstone or deletion certificate sufficient to identify the affected commitment, authority, reason, and frontier; conformance then applies to the surviving logical history.

\section{Observation, Rendering, and Attestation}

An observation is sampled information available to an observer or sensor. Rendering is the computation that produces a presentation from certified folds and observer conditions. Attestation is the distinct operation by which selected observed content becomes evidence that can constrain future admissibility.

\[
\operatorname{observe}: (M,Q,\omega,\rho)\rightarrow x \qquad [\text{no commit}],
\]
\[
\operatorname{attest}: (x,\text{witness},\text{rule},\phi)\rightarrow e \qquad [\text{candidate commit}].
\]

Rendering \MUST{} be side-effect-free with respect to $M$, $Q$, and $H$. Telemetry, access logs, billing records, or signed captures generated near a render call \MUST{} be modeled as separate operations. A cache miss, frame request, or additional viewer \MUSTNOT{} change world history unless an independently named policy declares that access itself is a world event.

An attestation \MUST{} identify the observed content or its content address, the witness or signing authority, the observation frontier if known, the rule under which it becomes constraining, and its verification status. Unverified material \MAY{} be committed as an allegation or pending attestation, but \MUSTNOT{} acquire verified standing merely by entering the log.

\begin{invariant}[Look/attest separation]
Repeated observation with identical certified inputs may change disposable caches and metrics, but cannot change authoritative folds. Only an admitted attestation event may make observed content historically constraining.
\end{invariant}

\section{Event Model}
\label{sec:event-model-section}

\subsection{Event envelope}

Each committed event \MUST{} include the following logical fields, though encoding and field names may differ:

\begin{longtable}{p{0.24\textwidth}p{0.65\textwidth}}
\toprule
Field & Contract \\
\midrule
Event identifier & Stable, globally unique or content-addressed identifier \\
World identifier & Namespace preventing cross-world replay \\
Kind and schema version & Typed payload discriminator and decoder version \\
Causal parents & Zero or more identifiers defining the event's immediate causal dependencies \\
Proposal reference & Identifier of the candidate being decided or transformed, when applicable \\
Disposition & Proposed, admitted, refused, deferred, repaired, attested, checkpointed, replay-discrepant, or a domain extension \\
Actor and authority & Principal initiating the candidate and authority admitting the event \\
Frontier & Causal frontier against which the gate evaluated the candidate \\
Payload & Canonically encoded domain data \\
Witness references & Bounded references to evidence not adequately represented by current folds \\
Versions & Gate, fold, rule, model, solver, and schema versions needed for interpretation \\
Integrity data & Payload hash and, where required, signatures or authenticated metadata \\
Time metadata & Optional physical time plus required logical or causal ordering data \\
\bottomrule
\end{longtable}

Physical timestamps \MUSTNOT{} be used as the sole causal-ordering mechanism. Canonical encoding \MUST{} define field ordering, number formats, string normalization, and treatment of unknown fields wherever hashes or signatures are compared.

\subsection{Candidate lifecycle}

A candidate begins outside authoritative history or as a committed \texttt{PROPOSED} event, according to domain policy. A gate evaluates
\[
A(c\mid M_{\phi},Q_{\phi},\phi,W)\in
\{\texttt{ADMIT},\texttt{REFUSE},\texttt{REPAIR},\texttt{DEFER}\}.
\]
The gate \MUST{} be versioned. Its decision \MUST{} name the evaluated frontier. Its witness set $W$ \MUST{} be bounded and content-addressed. Normal gate execution \SHOULD{} use certified folds; direct history lookup is reserved for evidence whose exact prior form matters.

\texttt{ADMIT} appends an admitted event and applies every relevant fold. \texttt{REFUSE} appends a refusal under the configured retention policy and normally leaves $M$ unchanged; it may update $Q$ only through an explicit rule. \texttt{REPAIR} appends an event that references both the defective or conflicting antecedent and the replacement disposition; it never overwrites either. \texttt{DEFER} appends a reason and blocking condition, changes neither $M$ nor ordinary permission standing by default, and creates an obligation governed by an explicit liveness policy.

A deferred candidate \MUST{} eventually be resolved, cancelled by an authorized event, or remain visibly pending. A timeout \MUSTNOT{} silently discard it. Resolution \MUST{} cite the deferral event as a causal parent.

\begin{invariant}[Disposition visibility]
Every gate invocation that the deployment policy defines as identity-bearing has exactly one durable disposition. Retries are correlated with the original candidate identifier and cannot masquerade as unrelated proposals.
\end{invariant}

\section{Semantic Folds}
\label{sec:semantic-folds}

A semantic fold is a deterministic, versioned reducer that converts a causally valid event history, or a checkpoint plus event suffix, into a queryable projection. The required authoritative folds are
\[
M_{\phi}=\fold(f_{\mathrm{phys}}^{v},M_0,H_{\phi}),\qquad
Q_{\phi}=\fold(f_{\mathrm{prov}}^{u},Q_0,H_{\phi}).
\]
Accounting, narrative, search, analytics, observer-index, and retention-summary folds \MAY{} be added. A derived fold does not become authoritative merely because it is cached.

Each fold definition \MUST{} declare its accepted event schemas, initial state, output schema, version, error behavior, ordering assumptions, treatment of unknown events, and migration policy. It \MUSTNOT{} consult undeclared ambient state such as current wall time, process-local randomness, network responses, mutable configuration, or unspecified floating-point behavior. Required external inputs must be captured in the event, checkpoint, or versioned execution profile.

If two events are concurrent, an implementation may apply them in either order only after demonstrating commutation for every authoritative fold they touch:
\[
f_k(f_k(s,e_i),e_j)=f_k(f_k(s,e_j),e_i)
\quad\text{for all affected }k.
\]
Commutation in $M$ alone does not imply commutation in $Q$. Conflicting events \MUST{} be serialized, refused, repaired, or deferred according to a deterministic policy.

\begin{invariant}[Prefix determinism]
Two conforming replicas using the same initial state, fold versions, canonical event content, and causally equivalent ordering of certified-independent events produce equivalent certified folds under the declared equality relation.
\end{invariant}

\begin{invariant}[Monotone history, selective folds]
Appending an event strictly advances logical history, but need not change every fold. A refusal or deferral normally lies in the kernel of the physical fold while remaining visible to provenance or pending-work folds.
\end{invariant}

\section{Optional Extension: Resonant or Mnemonic Fold}
\label{sec:resonant}

This section is non-normative for base conformance (Section~\ref{sec:conformance-levels}) and describes how a MEM|8-style continuously maintained associative memory can be added to the architecture without weakening any invariant already stated. Nothing in this section is required of a conforming implementation that does not adopt it.

The motivating question is whether a ``constant resonating pool'' -- a bounded, continuously updated memory whose present configuration reflects the accumulated effect of prior activity through resonance, interference, decay, and reinforcement -- is compatible with a history-primary architecture, given that rendering under Section~\ref{sec:appearance-section} \MUSTNOT{} scan $H_\phi$ on every frame. It is compatible, provided the pool is understood as a continuously maintained present disposition rather than as the complete history being rescanned, and provided its authority is bounded exactly as $M_\phi$'s and $Q_\phi$'s are bounded elsewhere in this specification.

\subsection{Definition and update rule}

An implementation \MAY{} define a resonant fold $\Psi_\phi$, distinct from the replay-profile vector $P$ of Section~\ref{sec:checkpoints-section} (the two are unrelated objects that happen to share initials in informal discussion; this specification reserves $P$ for the replay profile and $\Psi$ for the resonant fold to avoid the collision). The resonant fold is updated incrementally by each admitted event:
\[
\Psi_{t+1} = U(\Psi_t, E_{t+1}),
\]
where $U$ is a versioned, deterministic update rule combining the prior pool configuration with the newly admitted event's effect through whatever combination of resonance, interference, decay, and reinforcement the domain defines. Rendering \MAY{} then read the pool alongside the other certified folds:
\[
x_t = R(M_t, Q_t, \Psi_t, \omega_t).
\]
$\Psi_t$ therefore determines salience, expectation, associative readiness, and which continuations are easiest to recruit during presentation. It extends Section~\ref{sec:appearance-section}'s renderer signature by one argument; it does not replace $M_t$ or $Q_t$ in that signature, and an implementation that does not adopt this extension continues to render from $(M_t, Q_t, \omega_t)$ exactly as specified there.

\subsection{Fold family}

Adopting this extension produces a three-fold family over the same committed history, rather than the two-fold family of Section~\ref{sec:semantic-folds}:
\[
H_\phi \longrightarrow
\begin{cases}
M_\phi & \text{physical fold},\\
Q_\phi & \text{provenance fold},\\
\Psi_\phi & \text{resonant or mnemonic fold},
\end{cases}
\qquad
x_t = R(M_t, Q_t, \Psi_t, \omega_t).
\]
$\Psi_\phi$ \MUST{} be declared with the same fold obligations as any other fold under Section~\ref{sec:semantic-folds}: a declared accepted-event schema, initial state, output schema, version, and treatment of unknown events. Unlike $M$ and $Q$, $\Psi$ \MAY{} be lossy and non-invertible by design. A pool that remained constant in informational content while retaining arbitrary history exactly would not be a resonant fold in the intended sense; it would be a second copy of $H$ wearing different notation. The pool \MAY{} remain constant in capacity, dimensionality, and per-frame access cost only by being selective: some distinctions persist, some decay, and some become inseparably superposed. This lossiness is expected and unproblematic exactly because $\Psi$ is never the layer of record.

\subsection{Authority boundary}

\begin{invariant}[Resonant fold non-authority]
\label{inv:resonant-non-authority}
$\Psi_t$ \SHOULD{} influence rendering and \MAY{} influence heuristic prefiltering of candidates, but it \MUSTNOT{} be consulted by the admissibility gate $A$ of Section~\ref{sec:event-model-section} as a substitute for evidence drawn from $Q_\phi$ or from a bounded witness set $W$. A decision that would differ depending on the current configuration of $\Psi_t$ but not on $M_\phi$, $Q_\phi$, or $W$ \MUSTNOT{} be an admit, refuse, repair, or deferral disposition; it may only affect which candidates are surfaced for evaluation, never which are admitted.
\end{invariant}

This is the same boundary Invariant~1 (non-authority of appearance) already draws for $x_t$, extended to a second disposable-but-influential object. $\Psi_t$ is closer to appearance than to state in this respect: it is downstream of $H_\phi$ and causally shaped by it, exactly as required for the pool to be a faithful fold rather than an arbitrary cache, but it is not itself a source of admissibility, ownership, or identity. If reconstruction, audit, or dispute requires the precise past, the system \MUST{} consult $H_\phi$ through $M_\phi$ or $Q_\phi$; ordinary rendering or heuristic recall \MAY{} proceed from $\Psi_t$ alone. An implementation \MUSTNOT{} allow content that exists only in $\Psi_t$ -- and never became an admitted or attested event -- to acquire evidentiary standing merely because it is dynamically available for recall. The look/attest separation of Invariant~2 applies to the resonant fold without modification: resonance is a form of looking, however elaborate, until an attestation commits it.

\subsection{Relationship to the withdrawn replay-profile claim}

Section~\ref{sec:revision-guidance} withdraws, as an implementation claim, the suggestion that replay profiles $P=(p_b,p_n,p_s,p_c,p_p)$ form a useful join-semilattice, on the grounds that a componentwise maximum can combine guarantees achieved by different artifacts or executions and thereby describe no actual replay. The resonant fold does not reintroduce that error under a different name. $\Psi_t$ is not a summary of verification strength across runs; it is a single system's own accumulated disposition, updated by one deterministic rule along one causal history. It is a fold in the same sense $M_\phi$ and $Q_\phi$ are folds (Section~\ref{sec:semantic-folds}), not an aggregate statistic assembled after the fact across incomparable executions. The two constructions address different axes -- one asks how strongly a replay is verified, the other asks what a bounded present carries forward from an unbounded past -- and a conforming implementation \MUSTNOT{} conflate them by, for instance, treating a high-salience entry in $\Psi_t$ as evidence of a high replay grade for the event that produced it.

\subsection{Conformance}

Adopting this extension does not raise or lower an implementation's conformance level under Section~\ref{sec:conformance-levels}; the levels defined there are orthogonal to it. An implementation \MAY{} additionally claim \textbf{Resonant} compatibility by satisfying Invariant~\ref{inv:resonant-non-authority} and passing test \textbf{T13} below.

\begin{enumerate}[label=\textbf{T13}.,leftmargin=*]
\item \textbf{Resonant non-authority.} Populate $\Psi$ with content derived only from observation, never from an attested event. Confirm that every admit, refuse, repair, and deferral disposition produced by the gate is unchanged from a run with $\Psi$ held empty, and that no query against $M_\phi$ or $Q_\phi$ can retrieve that content as if it were certified.
\end{enumerate}

\section{Commit Protocol}

The minimal commit protocol has six phases. The implementation captures a candidate and its idempotency key; reads a certified $(M,Q,\phi)$; evaluates the versioned gate; constructs and durably appends one disposition event; applies or invalidates relevant fold caches; and publishes the new certified frontier. A crash before durable append leaves no commitment. A crash after durable append but before cache update is recovered by replay. Publication \MUSTNOT{} expose a certified frontier whose required authoritative folds are unavailable or fail verification.

Append and fold-cache update need not share one database transaction, but the history append is the commitment point. Cache records \MUST{} name their frontier and fold version. Consumers \MUST{} reject a cache whose frontier is neither equal to nor an explicitly acceptable ancestor of the requested frontier.

Idempotency is mandatory at the candidate and event levels. Reprocessing the same candidate \MUST{} return or reference its prior disposition rather than append a second logically independent admission. Content-identical candidates with distinct identifiers remain distinct proposals unless domain policy explicitly coalesces them.

\section{Checkpoints, Recovery, and Replay}
\label{sec:checkpoints-section}

A checkpoint is an accelerator bound to
\[
C=(\text{id},\phi,M_{\phi},Q_{\phi},V,\Xi,T,h).
\]
The tuple records its identifier, causal frontier, authoritative folds, software and schema versions, controlled stochastic inputs, comparison policies or tolerances, and integrity hashes. A checkpoint \MUST{} be derivable from or certified against history. It \MUSTNOT{} replace, truncate, or contradict history without a separately committed and auditable retention or migration event.

Recovery loads a verified checkpoint, verifies all declared bindings, and applies a causally valid suffix. Failure to verify a checkpoint or replay result \MUST{} append or externally preserve a \texttt{REPLAY\_DISCREPANT} record before a repaired checkpoint is promoted. The earlier artifact remains addressable, even if quarantined.

Replay claims are a vector, not a scalar:
\[
P=(p_b,p_n,p_s,p_c,p_p),
\]
covering bitwise, numerical, semantic, causal, and provenance recurrence. Each component is a predicate with declared inputs, environment, equality relation, tolerance, and evidence. A system \MUST{} report components independently. It \MUSTNOT{} infer semantic or provenance recurrence from bitwise or visual equality, nor claim a combined profile assembled from successes obtained by different runs unless it labels that profile as an aggregate rather than a single-run guarantee.

\section{Failure Cases and Required Responses}

\begin{longtable}{p{0.25\textwidth}p{0.31\textwidth}p{0.34\textwidth}}
\toprule
Failure & Detection & Required response \\
\midrule
Beautiful reset & Visual scene recurs but identifiers, repairs, or ownership do not & Reject persistence claim; restore from certified folds or record reconstruction as repair \\
Authoritative screenshot & Presentation is used as proof of complete or instantaneous global state & Require frontier and source-fold evidence; downgrade claim to presentation evidence \\
Invisible refusal & Rejected attempts leave no governed audit trace & Apply refusal retention/aggregation policy and update explicit provenance fold if authorized \\
Silent timeout & Deferred candidate disappears without disposition & Recover pending candidate and append cancellation, refusal, or resolution linked to the deferral \\
Repaired past & Divergent checkpoint is overwritten by a regenerated artifact & Preserve discrepancy and original identity; append repair and issue a new checkpoint \\
Poisoned ledger & Adversarial invalid proposals exhaust storage or reputation state & Apply versioned rate, aggregation, identity, and retention policies; certify summaries and their frontiers \\
Stale fold & Cache is served from a different or unknown frontier & Reject or label as stale; rebuild from a verified ancestor \\
Cross-fold race & Concurrent events commute in physics but conflict in ownership or permission & Treat as non-independent and serialize, refuse, repair, or defer \\
Schema drift & Old event gains new meaning under a new decoder & Retain old decoder semantics or commit an explicit migration producing a new fold version \\
Ambient nondeterminism & Replay changes because reducer consulted time, random state, locale, or network & Capture input, pin execution profile, append discrepancy, and re-certify \\
Observer-caused growth & Each render or sensor read appends world events & Separate telemetry from attestation and restore pure-render contract \\
Unauthorized attestation & Unsigned or ineligible observation alters admissibility & Mark as allegation/pending, reverse only through compensating events, and preserve incident lineage \\
Resonant leakage & Content present only in an optional resonant fold $\Psi$ (Section~\ref{sec:resonant}) is treated as certified evidence or changes a gate disposition & Reject the disposition, require an independent attestation, and audit the gate implementation for an undeclared $\Psi$ dependency \\
\bottomrule
\end{longtable}

\section{Security and Retention}

Append-only does not mean accept-all. Event admission \MUST{} authenticate the submitting principal where identity matters, authorize the claimed operation, validate schema and size, enforce namespace separation, and resist replay through candidate identifiers or nonces. Signatures attest to a signer and byte sequence; they do not establish truth. Truth or standing is assigned only through a declared gate and provenance rule.

Refusal retention is policy-sensitive because complete storage can enable denial of service while complete omission destroys evidence. A conforming policy \MUST{} state which refusals are retained individually, aggregated, sampled, or discarded; how aggregates bind to a frontier and rule version; and which security-relevant lineage cannot be compressed. Compression of history is itself a committed transformation with source range, algorithm version, output hash, custodian, and reversibility class.

Observer-specific rendering \MUST{} enforce $Q$ at the frontier from which $M$ is rendered. It \MUSTNOT{} render from one frontier while authorizing from an incompatible frontier unless the result is explicitly labeled stale and the disclosure risk is acceptable by policy.

\section{Rust Mapping}

The following mapping is illustrative but preserves the normative boundaries. Domain payloads should use canonical serialization rather than relying on compiler layout.

\begin{verbatim}
pub struct Frontier(pub Vec<EventId>);
pub struct EventId(pub [u8; 32]);
pub struct CandidateId(pub [u8; 16]);

pub struct EventEnvelope<P> {
    pub id: EventId,
    pub world: WorldId,
    pub schema: SchemaId,
    pub parents: Vec<EventId>,
    pub proposal: Option<CandidateId>,
    pub disposition: Disposition,
    pub actor: PrincipalId,
    pub authority: PrincipalId,
    pub evaluated_at: Frontier,
    pub payload: P,
    pub witnesses: Vec<WitnessRef>,
    pub versions: Versions,
    pub integrity: Integrity,
    pub physical_time: Option<Timestamp>,
}

pub enum Disposition {
    Proposed, Admitted, Refused(RefusalCode),
    Deferred { reason: DeferralCode, until: Condition },
    Repaired { antecedents: Vec<EventId> },
    Attested(Verification), Checkpointed(CheckpointId),
    ReplayDiscrepant(DiscrepancyId),
}

pub struct Certified<S> {
    pub value: S,
    pub frontier: Frontier,
    pub fold_version: FoldVersion,
    pub digest: Digest,
}

pub struct Provisional<S> {
    pub value: S,
    pub base: Frontier,
    pub execution: ExecutionProfile,
}

pub trait Fold<E, S> {
    type Error;
    fn version(&self) -> FoldVersion;
    fn apply(&self, state: &mut S, event: &E) -> Result<(), Self::Error>;
}

pub trait Gate<C, M, Q, P> {
    type Error;
    fn version(&self) -> GateVersion;
    fn decide(
        &self, candidate: &C, m: &Certified<M>, q: &Certified<Q>,
        witnesses: &[WitnessRef]
    ) -> Result<Decision<P>, Self::Error>;
}

pub trait Renderer<M, Q, O, X> {
    type Error;
    fn render(
        &self, m: &Certified<M>, q: &Certified<Q>, observer: &O,
        presentation: &PresentationProfile
    ) -> Result<X, Self::Error>;
}

pub trait History<E> {
    type Error;
    fn append(&self, expected: &Frontier, event: E)
        -> Result<Frontier, Self::Error>;
    fn suffix(&self, after: &Frontier) -> Result<Vec<E>, Self::Error>;
}
\end{verbatim}

Rust implementations \SHOULD{} make certified and provisional states distinct types, keep renderer access immutable, use exhaustive matching for dispositions, and return frontier conflicts as typed errors. The history implementation owns atomic append and idempotency. A fold does not append events; a gate does not mutate folds; orchestration is responsible for the commit protocol.

\section{Go Mapping}

Go cannot encode all boundaries in the type system, so interfaces and constructors must enforce them explicitly.

\begin{verbatim}
type Frontier struct { Heads []EventID }
type EventID [32]byte
type CandidateID [16]byte

type DispositionKind uint8
const (
    Proposed DispositionKind = iota
    Admitted
    Refused
    Deferred
    Repaired
    Attested
    Checkpointed
    ReplayDiscrepant
)

type EventEnvelope[P any] struct {
    ID          EventID
    World       WorldID
    Schema      SchemaID
    Parents     []EventID
    Proposal    *CandidateID
    Disposition Disposition
    Actor       PrincipalID
    Authority   PrincipalID
    EvaluatedAt Frontier
    Payload     P
    Witnesses   []WitnessRef
    Versions    Versions
    Integrity   Integrity
    PhysicalTime *time.Time
}

type Certified[S any] struct {
    Value       S
    Frontier    Frontier
    FoldVersion FoldVersion
    Digest      Digest
}

type Provisional[S any] struct {
    Value     S
    Base      Frontier
    Execution ExecutionProfile
}

type Fold[E, S any] interface {
    Version() FoldVersion
    Apply(context.Context, *S, E) error
}

type Gate[C, M, Q, P any] interface {
    Version() GateVersion
    Decide(context.Context, C, Certified[M], Certified[Q],
        []WitnessRef) (Decision[P], error)
}

type Renderer[M, Q, O, X any] interface {
    Render(context.Context, Certified[M], Certified[Q], O,
        PresentationProfile) (X, error)
}

type History[E any] interface {
    Append(context.Context, Frontier, E) (Frontier, error)
    Suffix(context.Context, Frontier) ([]E, error)
}
\end{verbatim}

Constructors \SHOULD{} defensively copy frontier heads and event slices. Implementations \MUSTNOT{} expose mutable maps or slices from certified values if callers could modify them without a commit. Disposition validation must reject impossible combinations, such as \texttt{Deferred} without a reason or \texttt{Repaired} without antecedents. As in Rust, the orchestrator composes history, gates, folds, and rendering; none silently performs another component's role.

\section{Conformance Test Suite}

An implementation claiming conformance \MUST{} publish which tests apply and provide machine-verifiable evidence for each supported deployment mode.

\begin{enumerate}[label=\textbf{T\arabic*.},leftmargin=*]
\item \textbf{Pure render.} Hash $M$, $Q$, and the history frontier before and after repeated rendering by multiple observers. All authoritative hashes and the frontier remain unchanged.
\item \textbf{Attestation boundary.} Observe identical content twice, attest only once, and verify that exactly one attestation lineage can constrain the gate.
\item \textbf{Fold reconstruction.} Rebuild $M$ and $Q$ from genesis and from each supported checkpoint. Results satisfy their declared equivalence relations at the same frontier.
\item \textbf{Idempotent retry.} Submit the same candidate identifier repeatedly across simulated crashes. At most one logical disposition is admitted.
\item \textbf{Refusal selectivity.} Refuse a candidate and confirm the physical fold is unchanged while the configured provenance or audit fold records exactly the policy-authorized effect.
\item \textbf{Deferral liveness.} Defer under a missing witness or partition, restore the condition, and verify that resolution cites the original deferral. A timeout cannot erase the candidate.
\item \textbf{Cross-fold concurrency.} Generate event pairs that commute in $M$ but conflict in $Q$. The system does not certify them as independent.
\item \textbf{Checkpoint corruption.} Alter a checkpoint byte. Recovery rejects it, preserves its identity, and records a discrepancy before promoting any repair.
\item \textbf{Version drift.} Replay under a new schema, fold, solver, or model version. The system either reproduces the declared profile or reports which component failed.
\item \textbf{Off-screen identity.} Remove an entity from all appearances, advance the world, and restore visibility. Identity and relevant properties derive from state and history, not regenerated resemblance.
\item \textbf{Observer isolation.} Render authorized and unauthorized views from the same frontier. Hidden content does not leak, and neither render changes permissions.
\item \textbf{History immutability.} Attempt in-place modification, deletion, reparenting, and reinterpretation. Each is rejected or represented as a new authorized event without changing prior event identity.
\end{enumerate}

\section{Conformance Levels}
\label{sec:conformance-levels}

\textbf{Layered} conformance requires explicit appearance, certified/provisional operative state, and pure rendering. \textbf{Historied} conformance additionally requires immutable typed events, deterministic reconstruction, dispositions, and checkpoint binding. \textbf{Provenance-preserving} conformance additionally requires a certified $Q$ fold and cross-fold conflict checks. \textbf{Distributed} conformance additionally requires causal frontiers, certified independence, partition-aware deferral or another explicit consistency policy, and replica convergence tests. An implementation \MUSTNOT{} claim a stronger level while omitting one of its obligations. An implementation \MAY{} separately claim \textbf{Resonant} compatibility (Section~\ref{sec:resonant}) by satisfying Invariant~\ref{inv:resonant-non-authority} and passing test T13; this label is orthogonal to the four levels above and neither raises nor lowers any of them.

These conformance levels describe engineering guarantees, not a metaphysical ranking. A short-lived game or visualization may rationally stop at the layered level. A collaborative scientific, economic, legal, or archival world will usually require provenance-preserving or distributed conformance.

\section{Revision Guidance for the Framework}
\label{sec:revision-guidance}

The next conceptual revision should preserve three sharp separations. First, ``observation'' should name disposable sampling, while ``attestation'' names the identity-bearing act that makes selected content constraining. Second, ``state'' should be split everywhere into certified operative state, provisional operative state, and provenance state; an unqualified $s_t$ obscures authority. Third, ``history-primary'' should not imply that rendering scans the log or that solver microsteps are events. History warrants folds at commitment boundaries; certified folds serve current queries; provisional state drives computation between boundaries.

The event vocabulary should treat proposal status and disposition as related but distinct. In some deployments a proposal is itself historically significant; in others it remains transient until the gate decides. Both are permitted if the retention policy is explicit. Deferral should remain first-class because it represents epistemic insufficiency rather than failure, and repair should always preserve its antecedents.

Replay should remain multidimensional. The earlier suggestion that replay profiles form a useful join-semilattice should be withdrawn as an implementation claim: a componentwise maximum can combine guarantees achieved by different artifacts or executions and thereby describe no actual replay. The safe object is an evidence-bearing vector attached to one declared run, plus separately labeled aggregates if needed.

Finally, the framework should distinguish invariants from aspirations. Pure rendering, frontier binding, deterministic folds, causal-parent integrity, idempotent commitment, and visible repair are enforceable invariants. Long-term identity, semantic equivalence, and the adequacy of a chosen event vocabulary are domain-relative claims that require declared tests and evidence.

\section{Compact Reference}

At runtime, render only from compatible certified folds; treat the result as presentation, not state. Evolve provisional physics from a named certified base. At a commitment boundary, evaluate a versioned gate against compatible $M$, $Q$, and frontier values. Append exactly one visible disposition using an idempotency key. Rebuild or advance every affected fold deterministically. Publish the new frontier only when required folds are available. Checkpoint the folds with their frontier, versions, inputs, tolerances, and hashes. On disagreement, append discrepancy and repair; never rewrite the past to resemble the desired present.

\end{document}
